% GENERATED FILE. Edit canonical lessons, then run npm run generate:books.
% canonical-source-hash: fnv1a64:e9105dbcbf50c3bb
% canonical-lessons: TE-C25-shubha-ratri, TE-S128-letter-bha

\chapter{Good Night}
\label{ch:te-good-night}
\hlchaptermodality{25}

\begin{chapteropening}
\textbf{By the end of this chapter:} \emph{I can wish someone good night in Telugu.}

Say \te{శుభ} \te{రాత్రి} — \textquotedblleft{}auspicious night\textquotedblright{} — pairing the \te{రాత్రి} already learned with śubha, while knowing many speakers code-switch to English here.
\end{chapteropening}

\section[śubha rātri]{\te{శుభ} \te{రాత్రి} (śubha rātri) — the standard translation, honestly checked against real speech}

\label{lesson:TE-C25-shubha-ratri}



\begin{quote}
You have both pieces already: \emph{rātri} from last lesson, and one more Sanskrit word. But this lesson's honest twist is about who actually says the result out loud.
\end{quote}

\begin{cousinweb}
\textbf{\te{శుభ} \te{రాత్రి}} (\textbf{śubha rātri}) — \textquotedblleft{}\textbf{good night}\textquotedblright{} — literally \textquotedblleft{}\textbf{auspicious night}.\textquotedblright{} Both pieces are Sanskrit \textbf{tatsama}: \textbf{\te{రాత్రి}} (\emph{rātri}, \textquotedblleft{}night,\textquotedblright{} already met) plus \textbf{\te{శుభ}} (\emph{śubha}, \textquotedblleft{}\textbf{auspicious},\textquotedblright{} from root \textbf{śubh}, \textquotedblleft{}to be beautiful,\textquotedblright{} ultimately PIE \emph{\textbf{*\'{k}ewb\textsuperscript{h}-}}).
\end{cousinweb}

\begin{culture}
Here's a genuinely interesting wrinkle, hedged appropriately: some sources describe modern Telugu speakers, especially casually, often \textbf{code- switching to English} \textquotedblleft{}\textbf{good night}\textquotedblright{} rather than saying \textbf{\te{శుభ} \te{రాత్రి}} aloud — treating it more as the standard written/dictionary translation than an everyday spoken phrase. Be careful with this claim, though: the evidence for it comes mainly from informal language-learning blogs and forum posts, not rigorous sociolinguistic sources, and the same \textquotedblleft{}code-switch to English for casual farewells\textquotedblright{} pattern plausibly shows up across other South Asian languages too — it isn't clearly something unique to Telugu, just something the available sources happened to mention for Telugu specifically. Either way, \emph{śubha rātri} itself is perfectly correct and understood Telugu, whatever its actual spoken frequency turns out to be.
\end{culture}

\subsection*{Your turn}
Say these aloud:

\begin{itemize}
\raggedright
  \item \textquotedblleft{}śubha rātri\textquotedblright{} — \textquotedblleft{}good night,\textquotedblright{} literally \textquotedblleft{}auspicious night\textquotedblright{}
  \item śubha's own root — śubh, \textquotedblleft{}to be beautiful\textquotedblright{} $\to$ \textquotedblleft{}auspicious\textquotedblright{}
  \item the honest twist — the standard translation, but many speakers code-switch to English instead
\end{itemize}

\begin{tcolorbox}[breakable,colback=teal!4,colframe=teal!35!black,title={Before you move on}]
What does \textbf{\te{శుభ} \te{రాత్రి}} literally mean? (\textquotedblleft{}\textbf{Auspicious night}.\textquotedblright{}) Is \emph{śubha rātri} actually the phrase most Telugu speakers say aloud for \textquotedblleft{}good night\textquotedblright{}? (\textbf{Often not} — many code-switch to English \textquotedblleft{}good night\textquotedblright{} in casual speech; \emph{śubha rātri} is the standard written/ formal translation.) Does this mean \emph{śubha rātri} is wrong or ungrammatical? (\textbf{No} — perfectly correct, just less commonly spoken than expected.)
\end{tcolorbox}
\section[bha]{\te{భ} — one character, met inside words you already say}

\label{lesson:TE-S128-letter-bha}



\begin{quote}
Before the new one: \te{జ} — what does it do? And one from further back: \te{ఢ}?

One character this time — and you have been saying it for pages without knowing which mark on the page it was.
\end{quote}

\subsection*{Script you'll notice: \te{భ}}
\textbf{\te{భ}} — \emph{bha}.

It is a \textbf{consonant}, and in this script a consonant is never bare: it comes with an \emph{a} already in it. So it is not \emph{bh}, it is \textbf{bha}.

Say \textbf{\te{బ}} (\emph{ba}), then say it again and let a puff of breath out on top of it. That second thing is \textbf{\te{భ}}. By now the pattern should arrive before the letter does: you have met \textbf{\te{క}/\te{ఖ}}, \textbf{\te{గ}/\te{ఘ}} and \textbf{\te{డ}/\te{ఢ}}, and this is the fourth pair of the same kind. The \emph{h} in the romanization \emph{bha} is that puff, not a separate sound.

You already say these, and every one of them has \te{భ} somewhere inside it:

\begin{itemize}
\raggedright
  \item \textbf{\te{శుభ} \te{రాత్రి}} \emph{śubha rātri} — good night, from the chapter before this one
  \item \textbf{\te{భాద్రపదం}} \emph{Bhādrapadaṁ} — the sixth of the twelve months
  \item \textbf{\te{ప్రారంభించు}} \emph{prārambhin̄cu} — to begin
\end{itemize}

\subsection*{Writing: \te{భ} — copy what you see}
Put your pen on \te{భ} and follow its line. Copy the shape you can see — slowly, and larger than it is printed.

\begin{quote}
This book does not yet tell you \textbf{where to start the character or which way to travel}. That is a real thing, taught with real variation from school to school, and it is not written down here until it can be written down with a source. Copying what is in front of you needs no such source, and it is how the shape gets into your hand in the meantime.
\end{quote}

\subsection*{Your turn}
\begin{itemize}
\raggedright
  \item \emph{Look:} at these words, and find \te{భ} in the ones that have it
\end{itemize}

\begin{quote}
\te{శుభ} \te{రాత్రి}  ·  \te{నమస్కారం}
\end{quote}

\begin{itemize}
\raggedright
  \item \emph{Say it:} \emph{ba}, then \emph{bha}, and hear the breath arrive
  \item \emph{Trace it:} \te{భ} three times, saying \emph{bha} as you finish each one
\end{itemize}

\begin{tcolorbox}[breakable,colback=teal!4,colframe=teal!35!black,title={Before you move on}]
Which character is this — \te{భ}? What does it carry that \te{బ} does not? (\emph{\textbf{bha}} — the same \emph{b} with a puff of breath after it.) Name one word you already say that contains it.

One more, from much earlier: \te{ద} — what does it do?
\end{tcolorbox}
